% P16-C4: Jin Tian 2004 conditional cancellation, with the corrected bridge.
\section{Tian's conditional-identification procedure}
\label{sec:p16-tian-conditional-identification}

Jin Tian's \emph{Identifying Conditional Causal Effects} gives a
four-phase procedure that identifies conditional effects by cancelling
an unknown factor shared by a numerator and denominator
\cite[Section~4, Eqs.~(25)--(30), Figure~3]{Tian2004ConditionalCausalEffects}.
Its control structure differs from IDC: it first attempts district
identification, then closes a set of interacting factors under the
variables being summed. The original paper leaves completeness open.
The later correction by Shpitser withdraws a claim of unsoundness based
on a notation misunderstanding and supplies a completeness argument
\cite{Shpitser2007ConditionalIdentificationAppendum}. The proofs below
retain Tian's procedure and derive its precise connection with IDC.

\begin{hbdefinition}{def:p16-tcc-data}{Tian's letters, positivity and the finite model domain}
Write \(V\) for the observed vertex set. The source's three
pairwise disjoint sets are \(T\) for actions, \(S\) for responses,
and \(C\) for conditioning variables, with values \(t,s,c\).
The target is
\[
 P_t(s\mid c)=P_t(s,c)/P_t(c),\qquad P_t(c)>0.
\]
Use the finite independent-latent-input class of
\cref{def:p16-latent-input-model},
and a compatible full observed law \(p(v)>0\) at every assignment.
Success soundness applies on any fixed finite nonempty alphabet family.
Graph-uniform completeness quantifies over every such family, every
possible common positive observed law, every two compatible models
sharing it, and all \(t,s,c\). Post-action conditioning masses
are positive by the indicator-removal support theorem.

Tian's Definition~1 explicitly uses positive observed distributions.
His Eqs.~(3)--(4) use normalized stochastic rows
\(k_i(v_i\mid v_{\operatorname{pa}(i)},u_i)\) and a latent
mixture. For the finite independent-input setting here, the following
bridge gives these rows the same intervention meaning as total equations.
No arbitrary dependent latent law is converted into independent
pair inputs by this statement.
\end{hbdefinition}

\begin{hblemma}{thm:p16-tcc-stochastic-bridge}{proved=annex:p16.tcc-stochastic-bridge}{Finite stochastic rows and total structural equations}
With finite independent latent inputs, every family of normalized
finite stochastic rows is represented by total acyclic equations
with additional independent private response tables. Observed and
all intervention laws agree with the source's mixture of products
and truncated products. Conversely deterministic equations give
those formulas using indicator rows. The declared observed mixed
graph is preserved.
\end{hblemma}

\begin{hbdefinition}{def:p16-tcc-inclusive-parents}{The source's inclusive parent scope}
For \(B\subseteq V\), define
\[
 \operatorname{Pa}^{+}_G(B)=B\cup\operatorname{pa}_G(B),
 \qquad \operatorname{pa}_G(B)=\bigcup_{i\in B}\operatorname{pa}_G(i).
\]
Tian writes \(\operatorname{Pa}(B)\) for this \emph{inclusive}
set, which contains the output vertices \(B\) themselves
\cite[Section~3, after Eq.~(5)]{Tian2004ConditionalCausalEffects}.
It bounds all arguments of \(Q[B](v)=P_{v_{V\setminus B}}(v_B)\).
Replacing it by proper parents in the conditional algorithm changes
the procedure. All Q kernels below retain their outside intervention
arguments, and \(\Sigma_B\) binds the indicated coordinates
of the common full assignment \(v\).
\end{hbdefinition}

\begin{hbmethod}{def:p16-tcc-component}{The source's component subroutine Identify}
On its component-call domain, input consists of nonempty connected
sets \(B\subseteq H\) and the full positive kernel \(Q[H]\),
with \(G[B],G[H]\) bidirected-connected. Put
\(A=\operatorname{an}_{G[H]}(B)\).
\begin{enumerate}
\item If \(A=B\), return \(\Sigma_{H\setminus B}Q[H]\).
\item Otherwise, if \(A=H\), return the retained failure record.
\item Otherwise \(B\subsetneq A\subsetneq H\). Compute
      \(Q[A]=\Sigma_{H\setminus A}Q[H]\), select the unique
      maximal district \(H'\) of \(G[A]\) containing \(B\),
      extract \(Q[H']\), and return
      \(\mathsf{Identify}(B,H',Q[H'])\).
\end{enumerate}
Here \(B\subseteq H'\subseteq A\subsetneq H\);
the first two inclusions need not be strict. Every use by the
conditional procedure has this component-call domain. An arbitrary
disconnected target is not silently passed to the assumption that
one containing district exists.
\end{hbmethod}

\begin{hblemma}{thm:p16-tcc-component-interface}{proved=annex:p16.tcc-component-interface}{The component interface and the conditional target decomposition}
The subroutine is the same component recursion already proved for
c-identify, with these renamed arguments; it terminates, preserves
positive full Q kernels, and is sound. Its failure criterion has
the earlier graph-uniform completeness meaning. The 2004 paper's
word ``unidentifiable'' in its subsequent analysis meant failure
of this procedure, since completeness was then left open.

Put \(D=\operatorname{an}_{G[V\setminus T]}(S\cup C)\),
\(F=D\setminus(S\cup C)\), and let \(D_1,\ldots,D_l\)
be the maximal districts of \(G[D]\). Then
\[
 P_t(s,c)=\Sigma_F\prod_{i=1}^l Q[D_i],
 \qquad \operatorname{Pa}^{+}_G(D_i)\subseteq D\cup T.
\]
Each \(D_i\) belongs to one full-graph maximal district \(H_i\),
whose kernel is identified by the topological conditional-factor formula.
\end{hblemma}

\begin{hbcorollary}{thm:p16-tcc-quick-criteria}{proved=annex:p16.tcc-quick-criteria}{The source's two preliminary conditional criteria}
For a singleton action vertex \(X\), the conditional effect is
identified if \(X\) has no bidirected path to any of its
directed children inside \(G[\operatorname{an}_G(S\cup C)]\).
Also, if no member of \(C\) is a descendant of any action
vertex, then \(P_t(C)=p(C)\), and conditional identification
of \(P_t(S\mid C)\) is equivalent to identification of
\(P_t(S,C)\). These are the preliminary cases in the source
before its general cancellation procedure.
\end{hbcorollary}

\begin{hbmethod}{def:p16-tcc-algorithm}{Algorithm 1: the original four phases}
The graph, full positive observed law and the value sets are input
along with \(T,S,C\). Empty products and sums over empty coordinate
sets have their usual unit and identity meanings.
\begin{enumerate}
\item \emph{Q-decomposition.} Find the full-graph maximal districts
      and compute their Q kernels from \(p\). Compute \(D,F\)
      and the districts \(D_i\) from the preceding lemma.
\item \emph{Component identification.} For each \(D_i\), call
      \(\mathsf{Identify}(D_i,H_i,Q[H_i])\). Put failed districts
      into \(\mathcal N\), and successful districts with their
      known Q expressions into \(\mathcal I\). If
      \(\mathcal N=\varnothing\), return
      \[
       \frac{\Sigma_F\prod_i Q[D_i]}
            {\Sigma_{S\cup F}\prod_i Q[D_i]}.
      \]
\item \emph{Closure of interacting factors.} Initialize
      \[
       H_0=\bigcup_{B\in\mathcal N}\operatorname{Pa}^{+}_G(B),
       \quad F_0=F\cap H_0,\quad F_1=F\setminus F_0,
       \quad\mathcal I_0=\varnothing,\quad\mathcal I_1=\mathcal I.
      \]
      Move every \(B\in\mathcal I_1\) with
      \(\operatorname{Pa}^{+}_G(B)\cap F_0\ne\varnothing\)
      into \(\mathcal I_0\). Put
      \(B_{\mathrm{new}}=F_1\cap
          \bigcup_{B\in\mathcal I_0}\operatorname{Pa}^{+}_G(B)\).
      If this is nonempty, move its vertices from \(F_1\) to
      \(F_0\) and repeat this phase's factor-movement step.
      Otherwise proceed to Phase 4.
\item \emph{Cancellation.} Put
      \(H'=\bigcup_{B\in\mathcal N\cup\mathcal I_0}
                      \operatorname{Pa}^{+}_G(B)\).
      If \(S\cap H'\ne\varnothing\), return the failure
      record, including the failed component calls and closure history.
      Otherwise return
      \[
       \frac{K(s,c,t)}{\sum_{s'}K(s',c,t)},
       \qquad K=\Sigma_{F_1}\prod_{B\in\mathcal I_1}Q[B].
      \]
\end{enumerate}
This is Tian's Figure~3 procedure, with his families \(N,I,I_0,I_1\)
printed in calligraphic type and his inclusive \(\operatorname{Pa}\)
made explicit. The numerator need not identify the original joint effect.
\end{hbmethod}

\begin{hbtheorem}{thm:p16-tcc-closure}{proved=annex:p16.tcc-closure}{Termination and least closure of the partition phase}
The closure phase terminates and produces the unique least pair
\((F_0,\mathcal I_0)\) satisfying the seed and propagation requirements.
Its result is independent of the order of factor movements. At termination,
\[
 F_0=F\cap H',\qquad
 \operatorname{Pa}^{+}_G(B)\cap F_0=\varnothing
       \quad(B\in\mathcal I_1).
\]
Thus \(F_0,F_1\) partition \(F\), and
\(\mathcal I_0,\mathcal I_1\) partition the successful districts.
\end{hbtheorem}

\begin{hbtheorem}{thm:p16-tcc-cancellation-soundness}{proved=annex:p16.tcc-cancellation-soundness}{The original cancellation proof and every denominator}
Every successful return is \(P_t(s\mid c)\). In Phase 4's
success case there is a positive, possibly unidentified factor
\(U(c,t)\) such that
\[
 P_t(s,c)=U(c,t)K(s,c,t),\qquad
 P_t(c)=U(c,t)\sum_{s'}K(s',c,t).
\]
The common factor does not use response coordinates and cancels
legitimately. All terms used to compute \(K\) are identified,
and the displayed denominator is positive.
\end{hbtheorem}

\begin{hbtheorem}{thm:p16-tcc-complexity}{proved=annex:p16.tcc-complexity}{Finite execution and polynomial graph work}
The whole procedure terminates. With \(n=|V|\), adjacency-matrix
graph operations and finite set operations can determine its graph
success/failure in \(O((n+1)^4)\) elementary operations. This
bound concerns the graphical procedure and compact formula construction;
it is not a polynomial bound in \(n\) for expanding arbitrary
probability tables or exact numeric arithmetic on unbounded alphabets.
\end{hbtheorem}

\subsection{The explicit bridge to IDC}

\begin{hbdefinition}{def:p16-tcc-incidence}{The factor-scope incidence graph}
Make one factor vertex \(f_i\) for each district \(D_i\),
and one variable vertex for each member of \(D\setminus C=S\cup F\).
Join \(f_i\) to \(v\) exactly when
\(v\in\operatorname{Pa}^{+}_G(D_i)\cap(D\setminus C)\).
These are undirected graph edges. Let \(\mathcal R\) be the
districts whose factor vertices lie in a connected component meeting
\(S\), and put \(D_{\mathcal R}=\bigcup_{B\in\mathcal R}B\).
Factor vertices and observed variable vertices are different kinds of
vertex even when a district is a singleton.
\end{hbdefinition}

\begin{hblemma}{thm:p16-tcc-moral-bridge}{proved=annex:p16.tcc-moral-bridge}{Factor-scope connectivity and the ancestral moral graph}
The response-connected factor blocks in the incidence graph correspond
exactly to response-connected blocks of the canonical latent-root
ancestral moral graph after the conditioning vertices \(C\) are deleted.
In particular a conditioning vertex \(c\in C\) is eligible
for IDC's Rule 2 exchange if and only if its district \(D_i\)
does not belong to \(\mathcal R\).
\end{hblemma}

\begin{hbtheorem}{thm:p16-tcc-idc-bridge}{proved=annex:p16.tcc-idc-bridge}{The closure decision and IDC's terminal joint problem}
The maximal IDC exchange set is
\(E=C\setminus D_{\mathcal R}\), and its remaining conditioning
set is \(C\cap D_{\mathcal R}\). The ancestor set of its
terminal joint query after acting on \(T\cup E\) is exactly
\(D_{\mathcal R}\), whose districts are the members of
\(\mathcal R\). Tian's conditional procedure fails exactly when
\(\mathcal N\cap\mathcal R\ne\varnothing\), which is
exactly failure of the terminal joint ID problem.
\end{hbtheorem}

\begin{hbcorollary}{thm:p16-tcc-completeness}{proved=annex:p16.tcc-completeness}{The corrected completeness conclusion}
Tian's conditional procedure is sound and complete on the stated
positive graph-uniform class, agrees with IDC on success versus failure,
and computes the same conditional effect whenever it succeeds.
This conclusion follows from the explicit bridge and the completed
IDC theorem, while the successful formula retains Tian's own cancellation
derivation. A failed record does not claim nonidentification at every
particular numerical law or on every separately prescribed small alphabet.
\end{hbcorollary}

\subsection{Original examples and the corrected comparison}

\begin{hbtheorem}{thm:p16-tcc-original-examples}{proved=annex:p16.tcc-original-examples}{The source's cancellation examples with all scopes retained}
For Tian's Figure~2 graph, the directed edges are
\(A\to B\to X\to Z\to Y\) and \(X\to W\to Y\),
and the bidirected edges are \(A\leftrightarrow X\),
\(A\leftrightarrow Y\), \(X\leftrightarrow W\).
Put \(H_1=\{A,X,W,Y\}\) and
\[
 q_1=p(y\mid z,w,x,b,a)p(w\mid x,b,a)p(x\mid b,a)p(a).
\]
Put \(q_Y=(\Sigma_{\{A\}}q_1)/(\Sigma_{\{A,Y\}}q_1)\).
The procedure identifies
\[
 P_x(y\mid w)=\sum_z p(z\mid x)q_Y(y;w,z),
\]
although the original joint on \(Y,W\) has an unidentified
\(Q[\{W\}]\) factor. Conditioning additionally on \(B\)
causes Phase 4 to fail through the \(\{A,Y\}\) district.

For the Figure~4 graph, use directed edges
\(X\to A\to B\to W\to Y\), \(X\to Z\to Y\),
and bidirected edges \(X\leftrightarrow A\),
\(X\leftrightarrow W\), \(X\leftrightarrow Y\),
\(W\leftrightarrow Y\). Put
\[
 q_2=p(y\mid z,w,b,a,x)p(w\mid b,a,x)p(a\mid x)p(x).
\]
For \(P_x(y\mid a)\), the only failed district is \(\{A\}\),
\(F_0=\varnothing\), and \(\mathcal I_0=\varnothing\).
Writing \(q_{WY}=\Sigma_{\{X,A\}}q_2\), its return is
\[
 \frac{\Sigma_{\{B,W,Z\}}\bigl(p(b\mid a)p(z\mid x)q_{WY}\bigr)}
      {\Sigma_{\{Y,B,W,Z\}}\bigl(p(b\mid a)p(z\mid x)q_{WY}\bigr)}.
\]
Dummy summed coordinates are distinct from retained action/evidence values.
The annex gives every component call and closure step for these cases.
\end{hbtheorem}

\begin{hbtheorem}{thm:p16-tcc-corrected-comparison}{proved=annex:p16.tcc-corrected-comparison}{The retracted UAI counterexample correctly returns failure}
Use the UAI comparison graph \(W\to Y\to Z\),
\(X\to Z\), \(X\leftrightarrow Z\), and query
\(P_x(w\mid z)\). Then \(D=\{W,Y,Z\}\),
\(F=\{Y\}\), \(\mathcal N=\{\{Z\}\}\), and
\(F_0=\{Y\}\). Because
\(Y\in\operatorname{Pa}^{+}(\{Y\})=\{W,Y\}\),
the Y factor enters \(\mathcal I_0\), whose scope reaches
the response \(W\). Phase 4 therefore returns failure.

For complete positive countermodels, let \(U,V,E_Y,E_Z\)
be independent, \(U,V\) fair and \(E_Y,E_Z\) of
one-probability \(1/4\). Use \(X=U,W=V,Y=W\oplus E_Y\)
and either \(Z=X\oplus U\oplus Y\oplus E_Z\) or
\(Z=Y\oplus E_Z\). Their common observed law is
\[
 p(x,w,y,z)=\tfrac14 q(y\oplus w)q(z\oplus y),
 \quad q(0)=3/4,\quad q(1)=1/4.
\]
The first conditional response-one probability is \(1/2\) at
either \(z\); the second is \(3/8\) at \(z=0\) and
\(5/8\) at \(z=1\). Both conditioning masses are \(1/2\).
This confirms the corrected failed decision and does not repeat the
withdrawn claim that Tian's procedure succeeds on this graph.
\end{hbtheorem}

\begin{proofboundary}{def:p16-tcc-scope-boundary}{The source model and support boundaries}
The original 2004 presentation uses a mixture of stochastic kernels
with upstream unobserved variables and positive observed distributions.
The finite independent-input class and its proved stochastic-row bridge
are the theorem domain here; no general latent projection is silently
assumed. Broader nonpositive execution requires the support licenses
and actual conditional denominators treated separately earlier.
The original open completeness questions are historical statements;
the completeness conclusion above is supplied by the explicit corrected
bridge. The 2007 addendum remains a distinct source, not a replacement
for reading and reconstructing Tian's original procedure.
\end{proofboundary}
